\documentclass[UTF8,a4paper,12pt]{ctexart}
\usepackage{geometry,amsmath,booktabs,multirow,hyperref}
\geometry{left=2.5cm,right=2.5cm,top=2.5cm,bottom=2.5cm}
\title{高性能芯片热管理系统问题五：敏感性与稳定性分析}
\author{数学建模项目组}\date{\today}
\begin{document}\maketitle

\section{问题与基准方案}
问题四得到的鲁棒设计为
\[
x^0=(x_1^0,x_2^0,x_3^0)=(0.20,3.5,6).
\]
问题二代理模型在该点的预测为
\[
(R_0,P_0,U_0)=(0.740179,0.123082,0.797950).
\]
本题考察加工误差和运行波动使 $x_1,x_2,x_3$ 偏离基准值时，三个性能指标如何变化，并评价方案是否稳定。

\section{扰动模型}
将连续扰动写为
\[
x_1=x_1^0(1+\delta_1),\quad x_2=x_2^0(1+\delta_2),\quad x_3=x_3^0+\delta_3.
\]
考虑加工误差范围 $\delta_1,\delta_2\in[-0.05,0.05]$，$\delta_3\in[-1,1]$。这里 $x_3$ 是排数，故采用绝对扰动一排更符合实际加工/设计误差。三个指标由问题二代理模型 $f_R,f_P,f_U$ 计算：
\[
R=f_R(x_1,x_2,x_3),\quad P=f_P(x_1,x_2,x_3),\quad U=f_U(x_1,x_2,x_3).
\]

运行工况方面，附件一给出入口质量流量 $\dot m_0=0.001\,\mathrm{kg/s}$。若流量存在 $\pm5\%$ 波动，则由流动阻力公式 $\Delta p\propto v^2\propto\dot m^2$，压降相对变化约为 $\pm10.25\%$；由冷却液携热热阻 $R_{fluid}=1/(\dot m c_p)$，其相对变化约为 $\mp5.26\%$。这给出工况波动的物理校核边界。入口温度变化主要表现为芯片绝对温度平移，在热阻定义 $R=(T_{max}-T_{in})/Q$ 下不会直接改变热阻，但会改变实际最高温度和可靠性裕度。

\section{单因素敏感性}
固定另外两个参数为基准值，逐一改变一个参数，得到：
\begin{table}[h]\centering\caption{基准方案的单因素预测}
\begin{tabular}{c|c|ccc}\toprule
扰动参数 & 水平 & $R$ & $P$ & $U$\\\midrule
\multirow{5}{*}{针肋宽度比} & 0 & 0.766333&0.106239&0.856799\\
&0.10&0.748057&0.108595&0.812882\\
&0.15&0.742818&0.114322&0.801792\\
&0.20&0.740179&0.123082&0.797950\\
&0.30&0.742698&0.149700&0.812003\\\midrule
\multirow{4}{*}{歧管深高比} &3.0&0.731605&0.144617&0.812034\\
&3.5&0.740179&0.123082&0.797950\\
&4.0&0.745297&0.107876&0.795029\\
&4.5&0.746960&0.099000&0.803271\\\midrule
\multirow{5}{*}{针肋排数} &2&0.750321&0.098422&0.789430\\
&4&0.744500&0.112114&0.792264\\
&6&0.740179&0.123082&0.797950\\
&8&0.737359&0.131325&0.806486\\
&10&0.736039&0.136845&0.817873\\
\bottomrule\end{tabular}\end{table}

为消除量纲影响，定义局部弹性系数
\[
E_{ij}=\frac{\partial y_j/y_j}{\partial x_i/x_i}.
\]
在基准点附近，脚本计算得到：
\begin{table}[h]\centering\caption{局部弹性系数}
\begin{tabular}{c|rrr}\toprule
参数 & $E_R$ & $E_P$ & $E_U$\\\midrule
$x_1$ 针肋宽度比 & -0.0072&0.3340&-0.0011\\
$x_2$ 歧管深高比 & 0.0647&-1.0448&-0.0746\\
$x_3$ 针肋排数 & -0.0145&0.2341&0.0267\\
\bottomrule\end{tabular}\end{table}

绝对值越大表示相对敏感性越强。基准点附近压降对歧管深高比最敏感，其次是针肋宽度比和排数；热阻与温度非均匀性的一阶弹性均较小，说明基准点位于相对平缓区域。

\section{联合蒙特卡洛分析}
独立抽取10000组扰动：$\delta_1,\delta_2\sim U[-0.05,0.05]$，$\delta_3\sim U[-1,1]$。得到的预测分布如下。
\begin{table}[h]\centering\caption{蒙特卡洛结果（10000次）}
\begin{tabular}{c|rrrr}\toprule
指标 & 均值 & 标准差 & 变异系数 & P5--P95\\\midrule
$R$&0.740191&0.001744&0.0024&0.737257--0.743090\\
$P$&0.123103&0.004809&0.0391&0.115275--0.131238\\
$U$&0.798333&0.002701&0.0034&0.793909--0.802914\\
\bottomrule\end{tabular}\end{table}

变异系数 $CV=\sigma/\mu$。可以看到，热阻和温度非均匀性的 $CV$ 分别仅为0.24\%和0.34\%，压降的 $CV$ 为3.91\%，仍处于较小范围。压降是本方案最需要关注的稳定性指标，主要由歧管深高比误差和针肋造成的局部阻力变化引起。

\section{稳定性判定与运行波动}
本文采用以下判据：
\begin{enumerate}
\item 输出变异系数 $CV<5\%$ 视为小扰动下稳定；
\item P5--P95 区间不跨越设计禁限值时，视为具有足够可靠裕度；
\item 若工况流量波动为 $\pm5\%$，按物理标度估计压降变化不超过 $\pm10.25\%$，冷却液携热热阻变化不超过 $\mp5.26\%$；该影响应在工程校核中叠加到代理模型输出上。
\end{enumerate}
按上述判据，基准鲁棒方案的几何扰动分析通过稳定性检验。其优势在于 $x_1=0.20$ 处于针肋强化换热与阻塞效应的折中区，$x_3=6$ 避开了排数过多导致温度不均匀性明显上升的区域；$x_2=3.5$ 虽对压降敏感，但在给定 $\pm5\%$ 加工误差下输出仍保持较小离散程度。

\section{结论}
在 $x_1,x_2$ 各 $\pm5\%$、$x_3$ 各 $\pm1$ 排的加工扰动下，推荐方案 $(0.20,3.5,6)$ 的预测均值约为 $(0.740191,0.123103,0.798333)$，热阻、压降和温度非均匀性变异系数分别为0.24\%、3.91\%和0.34\%。因此该方案对小范围几何误差具有较强稳定性。流量 $\pm5\%$ 波动时，应额外考虑压降约 $\pm10.25\%$、冷却液携热热阻约 $\mp5.26\%$ 的物理变化；如实际系统有严格压降上限，建议优先提高歧管加工精度并对流量进行闭环控制。

\end{document}
